\documentclass[aspectratio=169,t,xcolor=table]{beamer}
\usepackage[utf8]{inputenc}
\usepackage[brazil]{babel}
\usepackage{booktabs}
\usepackage{listings}

\usetheme{Ufg}
\usetikzlibrary{positioning, fit, backgrounds}

\setbeamertemplate{caption}[numbered]

% Configuração para o código JavaScript com coloração
\lstdefinelanguage{JavaScript}{
  keywords={const, let, async, await, throw, new, function, return, try, catch, if},
  keywordstyle=\color{UFGBlue}\bfseries,
  comment=[l]{//},
  commentstyle=\color{DarkGreen}\ttfamily,
  stringstyle=\color{DarkOrange}\ttfamily,
  morestring=[b]',
  morestring=[b]",
  morestring=[b]`
}
\lstset{
  basicstyle=\ttfamily\scriptsize,
  breaklines=true,
  showstringspaces=false
}

\AtBeginSection[]{%
    \setLayout{mainpoint}%
    \setBGColor{\currentsectioncolor}%
    \begin{frame}[noframenumbering]{}%
        \frametitle{\insertsectionhead}%
    \end{frame}%
    \setLayout{vertical}%
}

\setPrimaryColor{UFGBlue}
\setLogos{lib/logos/INF-compact-white.png}{lib/logos/INF-compact-white.png}

\begin{document}

\title[Seminário SD]{Service Discovery e Resiliência}
\subtitle{Da Arquitetura Local à AWS Cloud Native}
\author[Equipe]{
  Felipe Oliveira Carvalho, 
  João Pedro de Freitas Gonçalves, \\
  Eduardo Martins, 
  Wendel Souza, 
  Gabriel Machado Elias
}
\institute[UFG]{%
  Instituto de Informática\\
  Universidade Federal de Goiás
}
\date{2026}

\frame[noframenumbering]{\titlepage}

%---------------------------------------------------------
\setLayout{vertical}
\begin{frame}{Agenda}
    \tableofcontents
\end{frame}

%==========================================================
\section[UFGBlue]{Stack Tecnológica}

%--------------------------------------------------------
\setLayout{horizontal}
\begin{frame}{Tecnologias Utilizadas}

    \begin{columns}
        \column{0.5\textwidth}
        \textbf{Backend (Microsserviços)}
        \begin{itemize}
            \item Node.js 22 + TypeScript 5.7
            \item Fastify 5 (Alta performance)
            \item ES Modules
            \item JWT (Autenticação Stateless)
        \end{itemize}

        \vspace{0.5em}
        \textbf{Armazenamento}
        \begin{itemize}
            \item Redis 7 (Cache do feed em memória)
        \end{itemize}

        \column{0.5\textwidth}
        \textbf{Frontend (Web e API Gateway)}
        \begin{itemize}
            \item HTML5 + CSS3 (Glassmorphism)
            \item Vanilla JavaScript + Axios
        \end{itemize}

        \vspace{0.5em}
        \textbf{Infraestrutura e Orquestração}
        \begin{itemize}
            \item Docker e Docker Compose
            \item HashiCorp Consul (Service Discovery)
        \end{itemize}
    \end{columns}

\end{frame}

%==========================================================
\section[UFGBlue]{Arquitetura Atual (Docker + Consul)}

%--------------------------------------------------------
\setLayout{horizontal}
\begin{frame}{Arquitetura com HashiCorp Consul}

    \begin{columns}
        \column{0.55\textwidth}
        \begin{itemize}
            \item \textbf{Cérebro (Consul):} Registra e monitora.
            \item \textbf{Gateway (Frontend):} Descobre APIs.
            \item \textbf{Microsserviços:} Auth e Posts.
        \end{itemize}
        
        \vspace{0.5em}
        \begin{block}{Benefício}
            Portas dinâmicas garantem resiliência.
        \end{block}

        \column{0.45\textwidth}
        \begin{center}
        \begin{tikzpicture}[
            rect/.style={rectangle, rounded corners=3pt, minimum height=0.6cm, align=center, font=\scriptsize\bfseries, text=white},
            arr/.style={->, thick, color=gray!70}
        ]
            \node[rect, fill=UFGBlue, minimum width=3cm] (front) at (0, 3) {Frontend (Gateway)};
            \node[rect, fill=UFGBlue, minimum width=3cm] (consul) at (0, 1.5) {HashiCorp Consul};
            \node[rect, fill=UFGBlue, minimum width=1.4cm] (auth) at (-0.8, 0) {Auth};
            \node[rect, fill=UFGBlue, minimum width=1.4cm] (posts) at (0.8, 0) {Posts};

            \draw[arr] (front) -- node[right, font=\tiny] {Consulta} (consul);
            \draw[arr, dashed] (auth) -- (consul);
            \draw[arr, dashed] (posts) -- (consul);
        \end{tikzpicture}
        \end{center}
    \end{columns}

\end{frame}

%--------------------------------------------------------
\setLayout{vertical}
\begin{frame}[fragile]{Snippet: Anunciando Serviço no Consul}

    \begin{lstlisting}[language=JavaScript, caption=Registro Dinâmico (Backend)]
async function registerService({ name, host, port }) {
  const config = {
    ID: `${name}-${host}-${port}`,
    Name: name,
    Address: host,
    Port: port,
    Check: {
      HTTP: `http://${host}:${port}/health`,
      Interval: '10s'
    }
  };
  // Cadastra-se no Cérebro (Consul) ao ligar
  await axios.put(`http://consul:8500/v1/agent/service/register`, config);
}
    \end{lstlisting}

\end{frame}

%--------------------------------------------------------
\setLayout{vertical}
\begin{frame}[fragile]{Snippet: Descobrindo Serviços}

    \begin{lstlisting}[language=JavaScript, caption=Consulta Dinâmica (Frontend)]
async function discover(name) {
  // O Gateway pergunta ao Consul quem esta saudavel
  const url = `http://consul:8500/v1/health/service/${name}?passing=true`;
  const res = await axios.get(url);
  
  if (!res.data.length) throw new Error("Unavailable");

  // Balanceamento de carga basico (Random)
  const node = res.data[Math.floor(Math.random() * res.data.length)];
  return { address: node.Service.Address, port: node.Service.Port };
}
    \end{lstlisting}

\end{frame}

%==========================================================
\section[UFGBlue]{Resiliência: Graceful Degradation}

%--------------------------------------------------------
\setLayout{horizontal}
\begin{frame}{Graceful Degradation na Prática}

    \begin{columns}
        \column{0.5\textwidth}
        \textbf{Se o Auth Service cair:}
        \begin{itemize}
            \item Consul detecta a falha.
            \item Novos logins falham.
            \item \textbf{Feed continua rodando!}
        \end{itemize}

        \column{0.5\textwidth}
        \begin{block}{O Segredo: JWT Stateless}
            O Posts Service valida as sessões ativas localmente, sem depender de rede.
        \end{block}
    \end{columns}

\end{frame}

%--------------------------------------------------------
\setLayout{vertical}
\begin{frame}[fragile]{Snippet: Validação Stateless}

    \begin{lstlisting}[language=JavaScript, caption=Posts Service - Validação Desacoplada]
app.decorate("authenticate", async (req, reply) => {
  try {
    // Valida a assinatura matematicamente
    // Nao precisa consultar o Auth-Service!
    await req.jwtVerify();
  } catch (err) {
    reply.code(401).send({ error: "Token invalido" });
  }
});
    \end{lstlisting}

\end{frame}

%==========================================================
\section[UFGBlue]{Redes: Amazon VPC e Docker Network}

%--------------------------------------------------------
\setLayout{horizontal}
\begin{frame}{A Rede na Nuvem}

    \begin{columns}
        \column{0.5\textwidth}
        \textbf{VPC = Seu Docker Network Global}
        \begin{itemize}
            \item Assim como o Docker isola os contêineres do seu PC, a \textbf{VPC} isola os recursos na nuvem.
        \end{itemize}
        
        \vspace{0.5em}
        \textbf{Comparação Direta:}
        \begin{itemize}
            \item \texttt{docker network} $\rightarrow$ VPC Subnets
            \item \texttt{ports: 8080} $\rightarrow$ Security Groups
        \end{itemize}

        \column{0.5\textwidth}
        \begin{block}{Security Groups}
            Agem como firewalls. Permitem liberar apenas a porta do Frontend para a Internet, escondendo todo o restante do ecossistema.
        \end{block}
    \end{columns}

\end{frame}

%==========================================================
\section[UFGBlue]{A Etapa Intermediária (AWS EC2)}

%--------------------------------------------------------
\setLayout{horizontal}
\begin{frame}{Migração: Lift and Shift}

    \begin{columns}
        \column{0.55\textwidth}
        \textbf{Copiando o ambiente local:}
        \begin{itemize}
            \item O código não muda!
            \item Alugamos um servidor (EC2) e executamos o mesmo \texttt{docker-compose.yml}.
            \item \textbf{AWS CloudFront} entra como escudo HTTPS (CDN Global).
        \end{itemize}
        
        \vspace{0.5em}
        \begin{alertblock}{O Problema}
            A EC2 vira um gargalo e um Ponto Único de Falha (SPOF).
        \end{alertblock}

        \column{0.45\textwidth}
        \begin{center}
        \begin{tikzpicture}[
            rect/.style={rectangle, rounded corners=3pt, minimum height=0.5cm, align=center, font=\scriptsize\bfseries, text=white},
            arr/.style={->, thick, color=gray!70}
        ]
            \node[rect, fill=UFGBlue, minimum width=2.5cm] (cdn) at (0, 3) {CloudFront};
            \node[rect, fill=DarkGray, minimum width=2.5cm, minimum height=2.2cm] (ec2) at (0, 1) {};
            \node[font=\scriptsize\bfseries, text=white] at (0, 1.8) {Amazon EC2};
            
            \node[rect, fill=UFGBlue, minimum width=2cm] (front) at (0, 1.2) {Frontend};
            \node[rect, fill=UFGBlue, minimum width=0.9cm] (auth) at (-0.5, 0.4) {Auth};
            \node[rect, fill=UFGBlue, minimum width=0.9cm] (posts) at (0.5, 0.4) {Posts};

            \draw[arr] (cdn) -- (ec2);
            \draw[arr, dashed] (front) -- (auth);
            \draw[arr, dashed] (front) -- (posts);
        \end{tikzpicture}
        \end{center}
    \end{columns}

\end{frame}

%==========================================================
\section[UFGBlue]{Evolução: Arquitetura Cloud Native}

%--------------------------------------------------------
\setLayout{horizontal}
\begin{frame}{Quebrando a EC2 com Amazon ECS}

    \begin{columns}
        \column{0.55\textwidth}
        \textbf{Substituições na Nuvem:}
        \begin{itemize}
            \item EC2 $\rightarrow$ ECS Fargate (Serverless)
            \item Consul $\rightarrow$ AWS Cloud Map
            \item Redis $\rightarrow$ Amazon ElastiCache
        \end{itemize}
        
        \vspace{0.5em}
        \begin{block}{Escalonamento Real}
            Cada API escala de forma autônoma sem precisarmos provisionar hardware.
        \end{block}

        \column{0.45\textwidth}
        \begin{center}
        \begin{tikzpicture}[
            rect/.style={rectangle, rounded corners=3pt, minimum height=0.5cm, align=center, font=\scriptsize\bfseries, text=white},
            arr/.style={->, thick, color=gray!70}
        ]
            \node[rect, fill=UFGBlue, minimum width=3cm] (cdn) at (0, 3.5) {CloudFront};
            \node[rect, fill=UFGBlue, minimum width=3cm] (alb) at (0, 2.3) {ALB (Gateway)};
            \node[rect, fill=UFGBlue, minimum width=1.4cm] (auth) at (-0.8, 1) {ECS Auth};
            \node[rect, fill=UFGBlue, minimum width=1.4cm] (posts) at (0.8, 1) {ECS Posts};
            \node[rect, fill=UFGBlue, minimum width=3cm] (map) at (0, -0.2) {AWS Cloud Map};

            \draw[arr] (cdn) -- (alb);
            \draw[arr] (alb) -- (auth);
            \draw[arr] (alb) -- (posts);
            \draw[arr, dashed] (auth) -- (map);
            \draw[arr, dashed] (posts) -- (map);
        \end{tikzpicture}
        \end{center}
    \end{columns}

\end{frame}

%==========================================================
\section[UFGBlue]{Discovery Gerenciado (AWS Cloud Map)}

%--------------------------------------------------------
\setLayout{horizontal}
\begin{frame}{O Fim do Código de Descoberta}

    \begin{columns}
        \column{0.5\textwidth}
        \textbf{Service Discovery Serverless}
        \begin{itemize}
            \item O \textbf{Cloud Map} substitui o HashiCorp Consul.
            \item Em vez de \textit{API-based} (o Axios fazendo GET no Consul), passamos a usar \textit{DNS-based}.
        \end{itemize}
        
        \column{0.5\textwidth}
        \begin{block}{Magia da Rede (DNS)}
            O código deleta toda a complexidade do Consul. O Frontend simplesmente chama \texttt{http://auth.local}. O DNS da AWS descobre qual IP está vivo e roteia o pacote.
        \end{block}
    \end{columns}

\end{frame}

%--------------------------------------------------------
\setLayout{blank}
\setBGColor{UFGBlue}
\begin{frame}

    \centering
    \vspace{2cm}

    \textbf{\Huge Obrigado}

    \vspace{1.5cm}
    \begin{figure}
        \centering
        \includegraphics[height=3cm]{lib/logos/INF-full-white.png}
    \end{figure}

\end{frame}

%--------------------------------------------------------
\setLayout{titlepage}
\setBGColor{UFGBlue}
\titlepage

\end{document}
